% !TeX root = ../papers/nc-p3.tex
\section{階数一直線と二次曲線の配置の総次数}
\label{sec:rank-one-low-degree-arrangements}
\begin{theorem}\label{thm:rank-one-low-degree-arrangements}
被約な四辺有限の標準有限Frobenius窓の一般可逆な全階数反対根について、
階数一軌跡の全曲線成分の被約総次数は3以下である。
どの平面にも曲線成分は次数2までしか載らない。
\end{theorem}
\begin{proof}[平面二次部分と残余一点]
滑らかな二次曲線または交わる二直線からなる階数一部分 $D$ の入力平面を $h=0$、出力平面を $k=0$ とする。
滑らかな場合は平面二次の完全系で出力も平面二次である（次数2の整曲線）。
二直線の場合は各々が出力直線に同型に写り、共通点で交わる。平面上で
\[
             N|_{h=0}=aa^t+qB
\]
と書ける。$a$ は入出力平面を同型に結ぶ一次列、$q$ は $D$ の二次式、$B$ は定数対称行列である。
実際滑らかな二次曲線の同型は平面同型に延びる。二直線でも各線上の平方の線型表示は交点で符号を合わせれば
平面上の同じ一次列に延びる。$Bk\ne0$ は定理\ref{thm:opposite-column-divisor}による。
平面上の $D$ の外に階数一点があれば $B$ の階数は2以下であり、そこの $a$ は $B$ の像に入る。
これは二つの平方の非零な差の像が両方の列を含むからである。
$Bk\ne0$ なので像は出力平面のベクトル空間に含まれず、その交わりのベクトル次元は高々1である。
従ってそのような追加点は高々一点で、平面上に第三の直線も他の曲線もない。

列 $Nk$ を定数行操作すれば四式は $Q,h l_1,h l_2,h l_3$ となる。
$Q|_{h=0}=q$ と規格化でき、$l_i$ は独立である。$h$ と $Q$ が互いに素なのでこの基点の $h$ による残余イデアルは
$(Q,l_1,l_2,l_3)$ となり、残余長は高々1である。
特に $D$ から離れた階数一曲線の和 $F$ は総次数1以下である。
出力面 $k=0$ に別の曲線が載らないことは上の平面の議論を逆向きに適用すればよい。
逆向きも異なる反対根なので定理\ref{thm:all-finite-opposite-octics}により一般可逆である。
従って $F$ 上で出力の $k$ は次数 $\deg F$ の有限因子を切り、この全因子は $Nk=0$ に入る。
$h=0$ で $Nk=0$ の部分はちょうど $D$ なので、この因子は $h$ の外、長さ1以下の残余に入る。
\end{proof}
\begin{proof}[四本のねじれた直線]
三本の互いに素な直線は同じ母線族として一意の滑らかな二次曲面 $S$ に入る。
例えば二本を相補な二空間に取り、第三をその間の可逆線型写像のグラフにすればよい。
行列式は階数一各直線の一般点で少なくとも三重である。
$S$ への切断は $(8,8)$ なので同一族の三母線を三回ずつ含み得ず、$S$ が八次式を割る。
$S$ は重複因子でない（定理\ref{thm:rank-one-total-five}）。残余六次の $S$ への制限は
この三母線を二回ずつ含み、それらを除くと反対族の六母線だけである。
従ってこの三直線に加えて $S$ 上に階数一の整二次曲線はあり得ない。

四本が互いに素と仮定する。全て同一族なら反対族の各直線上の四つの交点の平方は全係数非零の関係を持つ。
その跡で出力四点も共線になり、入力直線の二次の包が出力直線の二次の包に写る。
従って $S$ 全体で階数二以下となり二次因子の重複に反する。
そうでなければ四つの三本組の二次曲面は相異なり、八次式はその四つの積となる。
各々滑らかで高々十四印を含むので重複度の和は $4\cdot14<60$ となり矛盾する。
\end{proof}
\begin{proof}[有理成分上の共通切断]
直線と二次曲線の被約和 $C$ が非退化で、総次数 $d$、連結成分数 $b$ とする。
各既約成分は有理で二つの超平面束 $L,L'$ の次数は同じである。その同型を各々で
二次座標から来る二乗の同型に合わせて取れる。各連結成分で成分の交わりグラフの全域木を選び、
辺ごとに選んだ一交点での同型の符号を根から順に合わせる。
すると二つの四座標空間は同じ次元 $d+b$ の空間に入る。実際各有理成分上の全切断の組に、
木の各辺で値を合わせる一条件を課した空間であり、新しい成分の一点評価は全射である。
従って共通切断は少なくとも $8-d-b$ 次元である。その積に対し二つの二乗の制限は
選んでいない交点を含めて全成分上で一致する。

$d+b\le6$ なら少なくとも二つを共有し、$C$ に二次式がない場合には
定理\ref{thm:rank-one-projected-quintics}の混合交代式の矛盾を得る。
$d+b\le5$ で二次式がただ一つの場合にも、三共通切断と $N_0=ww^t+qB$ の固定核による
定理\ref{thm:rank-one-quartic-normalization}の証明がそのまま適用できる。
これらの議論には $C$ の整性も全窓の印が $C$ に載ることも要らない。
ここで入力の非退化性だけを仮定してよい。出力が平面に入れば固定列の曲線次数上界2に反する
（以下使うのはいずれも $d\ge4$ である）。
\end{proof}
\begin{proof}[四次の二次二次完全交叉は可約でも生じない]
$d=4$ の連結な $C$ を仮定する。先の共通三切断により二次式が二つ以上必要である。
それらに共通平面があれば $C$ は平面と直線の和に入り、平面上に次数2までの上界から $d\le3$。
従って二つは固有に交わり、次数4の被約 $C$ は完全交叉そのものとなる。
二つの座標の二次の包は同一なので、出力も同じ型の完全交叉である。実際出力でも二次式は二つあり、
共通因子の可能性は同じ平面上界で除かれる。両側で線型正規、$\omega_C=\mathcal O_C$ である。
$L=L'$ なら座標は射影同値であり、定理\ref{thm:rank-one-quartic-normalization}の
$N=xx^t+A(q_1,q_2)$ と交代六方向の矛盾はこの被約完全交叉にも適用される。
よって $L\ne L'$。混合束の切断数は8である。各成分で逆束の次数が負なので随伴とRiemann--Rochによる。
従って曲線外印は八つ以上である。

二次式鉛筆における一群の曲線外印数の上界を再確認する。
既約な鉛筆元上の二印を結ぶ直線がその二次曲面に入らなければ、二印以外に二次曲面と交わらず
曲線とは素である。その直線を通る一般平面は滑らかな二次曲線を切り、そこに四つの動点と二印を持つ。
各成分上で平面鉛筆は非定数で一般切断は横断的である。
六つの相異なる点の平方は有理正規四次上で全係数非零の関係を持つ。従って入力でも六点は共面となる。
入力の対応平面は二印を常に含み、一定なら $C$ が退化するので動く。一般元は $C$ の成分を含まない。
二つの固定した入力印と所与の曲線成分を含む平面は高々一つである（その二印の直線が成分なら印が曲線上になる）。
従って一般元による四点切断が両側で等しく、$L=L'$ に反する。
これでその群は同一母線上であり、三印なら全母線が階数一なので例外は一群までとなる。
可約鉛筆元は異なる二平面である。重複平面は $C$ の被約性と非退化性に反する。
各平面は残りの二次式で切られた被約な二次部分を含む。実際その平面対と残りの式の完全交叉は $C$ そのもので、
両平面の交線を成分に持てばその一般点で重複し、被約性に反する。よって先の平面上界でその群は高々二印である。

したがって鉛筆値は四種類以上となる。先に得た三共通切断での小行列は
$N_0=ww^t+A(q_1,q_2)$ であり、その三行鉛筆はこれら全ての値で階数2以下である。
一般可逆なら三次行列式に四根が生じる。特異鉛筆の固定核、または固定核のない場合の零二行ブロックは
定理\ref{thm:rank-one-quartic-normalization}の混合評価の矛盾を与える。
これで連結四次部分を全て排除した。
\end{proof}
\begin{proof}[総次数の全分割]
総次数5の連結和なら $d+b=6$ なので二次曲面に入る。可約二次曲面には次数4までしか載らない。
滑らかな二次曲面での既約成分は二族の直線および $(1,1)$ の二次曲線である。
二次数 $(1,4)$ なら四本の同族直線、または同族三直線と一二次曲線を含み、前の二次曲面切断に反する。
$(2,3)$ なら大きい側の母線を一本除いて $(2,2)$ の被約部分が残る。これは二次式の切断で連結な完全交叉である。
二次錐の場合、直線は全て母線、整二次曲線は平面切断である。
五次には少なくとも一母線があり、その一本を除いた残りは二つの平面切断の和に書ける
（二母線の和は頂点を通る平面切断）。従って同じく被約な二次二次完全交叉が生じる。

非連結な和に次数2以上の連結成分があれば、それは整二次曲線または交わる二直線を含む。
最初の残余評価より、それから離れた全ての連結成分の総次数は1以下である。
よって非連結な五次なら連結四次があるか全てが互いに素な五直線となり、いずれも排除済みである。
非連結な四次なら残る可能性は連結三次とそれに素な一本の直線である。
これはどの二次曲面にも載らない。滑らかな二次曲面上の非平面連結三次は $(1,2)$ 型でどの直線とも正に交わる。
二次錐では母線同士が交わり平面二次切断も全母線に会う。
平面対なら連結三次は両平面に少なくとも一成分を持ち、同じ平面に入る別の直線と交わる。
しかし $d+b=6$ より二次式の不在は共通二切断の議論に反する。
総次数6以上は既に除かれている。これで全ての場合を尽くした。
\end{proof}
